\documentclass{article}
\usepackage{graphicx} % Required for inserting images
\usepackage{amsfonts} % For \mathbb
\usepackage{amsmath} 

\title{Cacoa-LM Methods}
\author{ }
\date{ }

\begin{document}

\maketitle

\paragraph{Overview.}
All analyses share one sample-level linear model, built from the sample metadata, and one set of sample relabelings drawn from that model. Each analysis differs only in the response it models: a gene's expression, a cell type's proportion, the distances between samples' expression profiles, or the same quantities measured in the neighbourhood of every single cell.

\section{Linear models and contrasts}

\paragraph{General linear model framework.}
Linear models were used to analyze different types of observable response features as a function of sample characteristics. For example, a differential expression test would correspond to modeling the average expression of a particular gene in a given cell type, conditioned on sample metadata such as disease status, batch identity, or age. 

Let $n$ denote the number of samples and $\mathbf{y} \in \mathbb{R}^n$ denote the response vector for a specific feature (e.g. log-transformed gene expression magnitude). The response was modeled using the general linear equation:
\begin{equation}
\mathbf{y} \;=\; \mathbf{F} \boldsymbol{\beta} + \boldsymbol{\epsilon},
\end{equation}
where $\mathbf{F} \in \mathbb{R}^{n \times p}$ is the design matrix constructed from the sample metadata, $\boldsymbol{\beta} \in \mathbb{R}^p$ is the vector of regression coefficients, and $\boldsymbol{\epsilon}$ represents the residual error term. A model is specified by the variable to be tested and, optionally, a formula naming the covariates to adjust for; without a formula the design holds the tested variable alone (plus any variables used to stratify the permutations, see below). A tested factor is coded without an intercept and placed first, so that its coefficients are the level means at the reference setting of the other covariates; other factors keep treatment coding and a tested numeric variable keeps the intercept. Terms without variation across the samples are pruned.

\paragraph{Contrasts and design partitioning.}
A test compares two settings of the metadata, for example two levels of a factor at the reference setting of the other covariates, or a step of one unit (one standard deviation by default) in a numeric variable. Each setting corresponds to a row of the design, and the contrast is their difference,
\begin{equation}
\mathbf{c} \;=\; \mathbf{f}_{\mathrm{target}} - \mathbf{f}_{\mathrm{reference}} \in \mathbb{R}^p ,
\end{equation}
so that transformed covariates and interactions are handled uniformly; a comparison may also be averaged over the levels of another factor. The estimated effect is $\mathbf{c}^\top \hat{\boldsymbol{\beta}}$ and the test statistic $S$ its studentized form, the estimate divided by its standard error.

For the permutation procedures the design is split along the contrast direction into one column carrying the tested effect and the remaining directions,
\begin{equation}
\mathbf{X} \;=\; \mathbf{F}\,\mathbf{c}\,/\,\mathbf{c}^\top\mathbf{c}, \qquad \mathbf{Z} \;=\; \mathbf{F}\,\mathbf{N},
\end{equation}
where $\mathbf{N}$ is an orthonormal basis of the null space of $\mathbf{c}$; the coefficient of $\mathbf{X}$ is the contrast itself, $\mathbf{Z}$ holds the nuisance directions (batch, other covariates), and $[\mathbf{X}\;\mathbf{Z}]$ spans the same space as $\mathbf{F}$.


\section{Statistics and background models}

\paragraph{Statistical Approach.}
Statistical significance was assessed by comparing the observed statistic $S_{obs}$ to a null distribution $\{S_{perm}^{(r)}\}_{r=1}^R$ generated by $R$ relabelings of the samples (one setting per analysis object, $R = 999$ by default, so that the smallest attainable p-value is $1/1000$). The p-value is the exceedance fraction with add-one smoothing,
\begin{equation}
P \;=\; \frac{1 + \sum_{r=1}^R \mathbb{I}(|S_{perm}^{(r)}| \ge |S_{obs}|)}{1 + R},
\end{equation}
where ties count as exceedances; one-sided tests drop the absolute values. The same relabelings are used for every response tested in an analysis (every gene, cell type or cell), so that statistics combined across responses are compared against a coherent null.

\paragraph{Randomization Strategies.}
The generation of the null distribution was adapted to the study design, supporting two user-selected modes:
\begin{itemize}
    \item \textbf{Block Randomization:} To preserve the stratification of the study, rows of the design matrix were permuted only within exchangeable blocks. By default, these blocks were defined by the intersection of all nuisance factors (all categorical variables present in $\mathbf{Z}$), ensuring that permutations do not confound variables such as batch or sex.
    \item \textbf{Freedman-Lane Procedure:} To control for nuisance variables $\mathbf{Z}$ while maintaining exact Type I error rates, the Freedman-Lane procedure is employed. The nuisance matrix was projected out of both the response and the core design. Let $\mathbf{H}_Z = \mathbf{Z}(\mathbf{Z}^\top \mathbf{Z})^{-1}\mathbf{Z}^\top$. The residualized response $\mathbf{R}_y = (\mathbf{I} - \mathbf{H}_Z)\mathbf{y}$ was computed, permuted within blocks to yield $\tilde{\mathbf{y}}^*$, and regressed against the residualized core design $\mathbf{R}_X = (\mathbf{I} - \mathbf{H}_Z)\mathbf{X}$ to obtain null statistics.
\end{itemize}
A relabeling is a permutation of the samples. For the pairwise distance models of Section 3 the distances are left in place: under block randomization the permutation is applied to the rows of the design, under the Freedman-Lane procedure to the residual part of the samples' embedding, which preserves the dependence between pairs that share a sample.

\paragraph{Handling of Missing Data.}
Missing values (NAs) in the response vector were handled via two distinct user-selected modes:
\begin{itemize}
    \item \textbf{Drop Mode:} Observations containing NAs were explicitly excluded from the model fit for the specific feature. The design matrix was subset to the valid rows.
    \item \textbf{Impute-Weak Mode:} NAs were imputed with the mean of the observed values (or zero) and assigned a negligible weight ($w \approx 10^{-4}$) in the resulting (weighted) linear model fit. This approach maintains constant vector dimensions, while minimizing the influence of missing data on the estimated coefficients.
\end{itemize}
In the per-cell analyses (Section 4) a sample that has too few cells in a cell's neighbourhood is dropped for that cell only, as in drop mode; a cell is left untested under a relabeling that leaves fewer than two samples in a compared group.

\paragraph{Robust estimation.}
To mitigate the effect of outliers, coefficients were estimated using robust methods when specified. The ``Winsorized'' approach clamps the residuals $r_i$ at $k$ robust standard deviations, $\tilde{r}_i = \text{sign}(r_i) \cdot \min(|r_i|, k \cdot \text{MAD}(\mathbf{r}))$, and refits the model to the clamped responses; Huber iteratively reweighted least squares down-weights the observations with large residuals instead, with the same tuning constant $k = 1.345$. Under either scheme the weights are re-estimated under every relabeling.

\paragraph{Inference across many units.}
Several analyses test thousands of units at once (every cell of a map, every cell type of a table), with one relabeling set shared by all of them. Let $s_i$ be the statistic of unit $i$, $s_i^{(r)}$ its value under relabeling $r$, $m$ the number of units, and let the statistics be on a common scale across units (a studentized or standardized statistic). Two quantities are computed from the pooled null, the collection of all $m R$ permuted values:
\begin{itemize}
    \item \textbf{Pooled p-value.} The rank of a unit's statistic among all permuted values of all units,
    \begin{equation}
    P_i \;=\; \frac{1 + \sum_{j=1}^{m}\sum_{r=1}^{R} \mathbb{I}\big(s_j^{(r)} \ge s_i\big)}{1 + mR},
    \end{equation}
    which has no per-unit floor at $1/(R+1)$ and is calibrated whenever the units' null distributions are alike.
    \item \textbf{Pooled false-discovery rate.} For a threshold $t$, the expected number of null units above it, estimated from the relabelings, against the observed number:
    \begin{equation}
    \mathrm{FDR}(t) \;=\; \frac{\frac{1}{R}\sum_{r=1}^{R} \#\{j : s_j^{(r)} \ge t\}}{\#\{i : s_i \ge t\}},
    \qquad
    q_i \;=\; \min_{t \le s_i} \mathrm{FDR}(t),
    \end{equation}
    the q-value of a unit being the smallest rate of any rejection set that contains it.
\end{itemize}
For units that sit on a graph (cells), the statistic of a unit is replaced, before these two steps, by its mean over the unit and its graph neighbours, in the observed data and under every relabeling alike,
\begin{equation}
\bar{s}_i \;=\; \frac{1}{|\mathcal{N}_i|}\sum_{j \in \mathcal{N}_i} s_j ,
\end{equation}
which borrows strength from the neighbours while remaining an exact permutation test. Group-level summaries use the same relabelings: the mean statistic of a cell type's cells is compared with the distribution of the same mean under the relabelings, and the resulting p-values are adjusted across cell types by the Benjamini--Hochberg procedure. Per-unit p-values are reported as z-scores, the signed normal quantile of the p-value ($\mathrm{sign}(\text{effect}) \cdot \Phi^{-1}(1 - P_i/2)$ for a two-sided test), and adjusted z-scores are the same transform of the q-values.

\section{Pairwise distance models}

\paragraph{Pairwise distance analysis context.}
The expression state of a cell type in a sample is a vector over thousands of genes, and a condition may change it in many genes at once, each by a little. To ask whether the state changes between conditions without choosing genes, the analysis works with the distances between samples: for each cell type, expression is summed over the cells of each sample (pseudobulk), log-transformed, and a distance is computed between every pair of samples (the correlation distance by default). The questions are then whether samples of different conditions lie farther apart than samples of the same condition, and whether the samples of one condition scatter more than those of the other.

A direct regression of the pair distances on pair-level covariates answers these questions poorly: pairs that share a sample are not independent, the mean distance between two groups mixes the separation of the groups with the scatter inside each of them, and covariates are awkward to express at the level of pairs. The approach taken here instead treats the distance matrix as a description of where the samples sit relative to one another. A distance matrix of this kind can always be realized by points in a Euclidean space (the space that classical multidimensional scaling would draw), and in that space the questions become ordinary questions about samples rather than pairs: where does a sample's position depend on its metadata, and how far does it scatter around that position. The sample-level design, contrasts and relabelings of Sections 1 and 2 then apply unchanged.

\paragraph{Embedding of the samples.}
The positions are never computed. Every quantity needed below depends on the positions only through their inner products, and those are obtained from the $n \times n$ distance matrix $\mathbf{D}$ by double centring its squared entries,
\begin{equation}
\mathbf{G} \;=\; -\tfrac{1}{2}\,\mathbf{J}\,\mathbf{D}^{2}\,\mathbf{J}, \qquad \mathbf{J} = \mathbf{I} - \tfrac{1}{n}\mathbf{1}\mathbf{1}^\top ,
\end{equation}
so that $G_{ij}$ is the inner product of the positions of samples $i$ and $j$ about their common centre. This is also what makes the method indifferent to the choice of distance, as long as it is of Euclidean type.

\paragraph{Location and dispersion models.}
Two models are fitted on this space, one for each question. The location model places the position $\mathbf{y}_i$ of sample $i$ according to its row $\mathbf{x}_i$ of the sample-level design, $\mathbf{y}_i = \mathbf{B}^\top \mathbf{x}_i + \mathbf{e}_i$. Its least-squares fit depends on the data only through
\begin{equation}
\mathbf{M} \;=\; \mathbf{A}\,\mathbf{G}\,\mathbf{A}^\top, \qquad \mathbf{A} = (\mathbf{X}^\top\mathbf{X})^{-}\mathbf{X}^\top,
\end{equation}
the matrix of inner products between the fitted effects of the design columns ($\mathbf{M} = \mathbf{B}\mathbf{B}^\top$ for Euclidean data); the squared distance between the fitted positions of any two settings of the metadata is read off $\mathbf{M}$. The dispersion model describes how far each sample scatters around its fitted position, $s_i = \mathbb{E}\|\mathbf{e}_i\|^2 = \mathbf{z}_i^\top \boldsymbol{\gamma}$, where $\mathbf{z}_i$ is a row of a dispersion design (by default the compared groups, so that each group has its own scatter). It is fitted from the squared residual distances $r_i = (\mathbf{R}\mathbf{G}\mathbf{R})_{ii}$, $\mathbf{R} = \mathbf{I} - \mathbf{X}\mathbf{A}$, through the identity $\mathbb{E}[\mathbf{r}] = (\mathbf{R} \circ \mathbf{R})\,\mathbf{Z}\boldsymbol{\gamma}$, which holds whatever the pattern of unequal scatter. The dispersions are needed for a second reason: the fitted separation of two groups contains their scatter (two groups with the same centre and scatter $s$ still show an expected squared distance $s/n_T + s/n_R$ between their fitted positions), so the location effect is corrected, $\mathbf{M} \leftarrow \mathbf{M} - \mathbf{A}\,\mathrm{diag}(\mathbf{s})\,\mathbf{A}^\top$, and a true shift of zero is reported as zero on average.

\paragraph{Effects.}
For a contrast between two endpoints with design rows $\mathbf{x}_T$ and $\mathbf{x}_R$ (target and reference) and $\mathbf{c} = \mathbf{x}_T - \mathbf{x}_R$, three effects are reported:
\begin{equation}
\mathrm{shift} \;=\; \mathbf{c}^\top \mathbf{M}\,\mathbf{c}, \qquad
\mathrm{var} \;=\; 2\,(s_T - s_R), \qquad
\mathrm{total} \;=\; \mathrm{shift} + \tfrac{1}{2}\,\mathrm{var}.
\end{equation}
The shift is the squared distance between the two endpoints' positions, free of their scatter; the variance effect compares the scatter of the two groups; the total is the expected excess of a between-group distance over a within-reference distance. In a one-factor design they equal the classical quantities computed from mean pair distances, $m_{RT} - \tfrac{1}{2}(m_{RR} + m_{TT})$, $m_{TT} - m_{RR}$ and $m_{RT} - m_{RR}$. For reporting they are normalized by the scatter: $\mathrm{shift}/(s_T + s_R)$, $\log_2(s_T/s_R)$ and $\mathrm{total}/(2 s_R)$.

\paragraph{Tests.}
The shift of a contrast is tested with the permutation F statistic of the distance-based linear model,
\begin{equation}
F \;=\; \frac{\mathbf{a}^\top \mathbf{G}\,\mathbf{a}\;/\;\mathbf{c}^\top(\mathbf{X}^\top\mathbf{X})^{-}\mathbf{c}}{\mathrm{tr}(\mathbf{R}\mathbf{G})\;/\;(n - q)}, \qquad \mathbf{a} = \mathbf{A}^\top \mathbf{c},
\end{equation}
with $q$ the rank of the design, under the relabelings of Section 2. A multi-level term is tested by the same ratio with the term's sum of squares, $\mathrm{tr}((\mathbf{H}_{full} - \mathbf{H}_{reduced})\mathbf{G})$, in the numerator. The variance effect is tested by regressing the square roots of the leverage-corrected per-sample dispersions, $\sqrt{r_i/(1 - h_i)}$, on the dispersion design. The influence of single samples is assessed by refitting the effects without each sample in turn.

\section{Cluster-free maps}

\paragraph{The unit.}
Density and expression changes are also mapped at the level of single cells, without clusters. Each cell is tested in its own neighbourhood on the cell graph: the cell, all of its graph neighbours and, if those are fewer than $k - 1$, the two-hop neighbours with the most shared neighbours, up to $k$ cells. The size is set from the data as the larger of 100 cells and 10 cells per sample, capped at 1\% of the cells but never below 5 cells per sample. The neighbourhood is never reduced to the cell's strongest edges: when expression differs between the compared conditions, the strongest edges join cells of the same condition, and a unit built from them would hide the very composition change it is meant to test. A diffusion kernel on the graph is available as an alternative unit, with its width stated as the equivalent number of equally weighted cells.

\paragraph{Density map.}
For cell $i$ and sample $s$ let $c_{is}$ be the number of the sample's cells in the neighbourhood and $n_s$ the sample's size. The response
\begin{equation}
y_{is} \;=\; \log_2\!\left(\frac{c_{is} + 1/2}{n_s + 1}\cdot 10^{4}\right)
\end{equation}
is modeled across samples with the linear model and contrast of Section 1; the reported effect is the $\log_2$ fold change of the neighbourhood's share of a sample between the contrast endpoints, and the statistic is the studentized contrast $t$.

\paragraph{Expression-shift map.}
For each cell, the neighbourhood's cells are summed per sample into expression profiles, the distances between samples are computed, and the location model of Section 3 is fitted and tested exactly as for a cell type, with the shift as the effect and its $F$ as the statistic.

\paragraph{Inference.}
Both maps are adjusted as described in Section 2: the per-cell statistic is averaged over the cell and its graph neighbours, pooled p-values and pooled false-discovery rates are computed over all cells, and the maps show the corresponding z-scores. The pooled statistic is the studentized $t$ for the density map and the permutation-standardized $F$ for the shift map. Per-cell-type summaries use the group-level test of Section 2.

\section{Covariate screen}

Before a model is fixed, each metadata variable can be screened for association with a cell type's expression state. For every cell type and variable, the term test of Section 3 is applied to the variable alone and to the variable added to the other variables, giving marginal and partial effects (as fractions of distance variance explained) with Freedman--Lane permutation p-values on the same distance matrices; the association of the variable with the per-sample dispersions is tested alongside. P-values are adjusted over the whole grid of cell types and variables, and a global p-value per variable is obtained from the maximum statistic across cell types under the shared relabelings.

\section{CoDA Linear Model and Statistics}

\paragraph{Compositional linear model in ILR space and cell-type loadings.}
Let $n$ denote the number of samples and $D$ the number of cell types (parts). For each sample $i \in \{1,\dots,n\}$ we observe a composition (cell-type frequency vector) $x_i = (x_{i1},\dots,x_{iD})$ with $x_{ic} > 0$ and $\sum_{c=1}^D x_{ic}=1$. Zero counts were replaced by a pseudocount of 0.1 before forming the proportions. Let $K = D-1$ and $\Psi \in \mathbb{R}^{D \times K}$ be an ILR basis matrix with orthonormal columns, i.e.\ $\Psi^\top \Psi = I_K$, spanning the clr subspace.

We compute ILR coordinates by
\begin{equation}
Z \;=\; \log(X)\,\Psi \in \mathbb{R}^{n \times K},
\end{equation}
where $X \in \mathbb{R}^{n \times D}$ stacks the $x_i$ row-wise. We then fit a multivariate linear model in ILR space,
\begin{equation}
Z \;=\; F B + E,
\end{equation}
where $F \in \mathbb{R}^{n \times p}$ is the design matrix (with $p$ columns, e.g. intercept, group indicators, other covariates), $B \in \mathbb{R}^{p \times K}$ is the matrix of regression coefficients, and $E$ are the residuals.

\paragraph{Coefficients, contrasts, and mapping back to cell types.}
For a specific model coefficient (column of $F$ design matrix) $j$, the ILR-space effect is the $j$-th row of $B$,
\begin{equation}
\beta^{(j)} \;=\; B_{j,\cdot} \in \mathbb{R}^{K}.
\end{equation}
For a user-defined contrast vector $\theta \in \mathbb{R}^{p}$ (defined on the same coefficient ordering as $B$), the ILR-space contrast effect is:
\begin{equation}
\beta^{(\theta)} \;=\; \theta^\top B \in \mathbb{R}^{K}.
\end{equation}

To obtain effects for a given covariate or contrast at the level of individual cell types, this ILR-space
coefficient vector is mapped back to the cell-type (clr) space using the
ILR basis matrix $\Psi \in \mathbf{R}^{D \times (D-1)}$,
\begin{equation}
\ell^{(j)} \;=\; \Psi\,\beta^{(j)} \in \mathbf{R}^{D}.
\end{equation}
Each element $\ell^{(j)}_c$ represents the loading for cell type $c$
associated with covariate $j$. By construction, the loadings satisfy the
zero-sum constraint
\begin{equation}
\sum_{c=1}^{D} \ell^{(j)}_c \;=\; 0,
\end{equation}
and are therefore interpretable only in relative terms, i.e.\ as changes
in log-ratios of cell-type proportions rather than absolute changes.


\paragraph{Interpretation of $\ell$ as log-fold changes of proportions.}
The loadings $\ell$ live in clr (log-ratio) space and are interpretable only in relative terms. For any two cell types $a,b$,
\begin{equation}
\ell_a - \ell_b
\;=\;
\Delta \log\!\left(\frac{x_a}{x_b}\right),
\end{equation}
where $\Delta$ denotes the expected change associated with a one-unit change in the covariate (or moving along the contrast). Hence, differences of loadings correspond to changes in log-ratios of cell-type proportions. Equivalently, each $\ell_c$ can be viewed as the expected change in $\log(x_c)$ relative to the average change across all parts (the geometric mean change), i.e.
\begin{equation}
\ell_c
\;\approx\;
\Delta \log(x_c)
\;-\;
\frac{1}{D}\sum_{m=1}^{D}\Delta \log(x_m).
\end{equation}
For presentation as log$_2$-fold changes, we report $\ell_c/\log(2)$, so that a value of $+1$ corresponds to an approximate two-fold increase of cell type $c$ relative to the chosen baseline (e.g.\ a reference set), and $-1$ to an approximate two-fold decrease.



\paragraph{Statistical assessment of covariate and contrast effects.}
Statistical significance of compositional effects was assessed using
sample randomization, either via block permutations or the
Freedman-Lane procedure, applied consistently to the fitted linear
model in ILR space. For each covariate or user-defined contrast, both
global (composition-level) and cell-type–specific statistics were
computed and evaluated against corresponding null distributions
obtained under randomization.

For each covariate $j$ (with ILR-space effect $\beta^{(j)}$) and
for each user-defined contrast $\theta$ (with effect
$\beta^{(\theta)}$), a global composition-level statistic was computed
based on the fraction of ILR variance explained by the corresponding
fitted component. For each covariate $j$ (with ILR-space effect
$\beta^{(j)} = B_{j,\cdot}$) we define $v = F_{\cdot j} \in \mathbf{R}^{n}$
as the $j$-th column of the design matrix; for a user-defined contrast
$\theta \in \mathbf{R}^{p}$ (with effect $\beta^{(\theta)} = \theta^\top B$)
we define $v = F\theta \in \mathbf{R}^{n}$. The fitted ILR contribution
of the covariate/contrast is then
$\widehat{Z}_{\mathrm{effect}} = v\,\beta^\top$, and the global
composition-level statistic is the fraction of ILR variance explained,
\[
T \;=\; \frac{\|\widehat{Z}_{\mathrm{effect}}\|_F^2}{\|Z\|_F^2}.
\]
Empirical p-values for $T$ were obtained by comparing the observed value
to its null distribution computed from analogous quantities derived
from randomized fits. 

To assess cell-type–specific contributions, the ILR-space effect
$\beta$ was mapped back to the cell-type (clr) space as
$\ell = \Psi\,\beta$. For each cell type $c$, the loading $\ell_c$ was
used as a cell-type–level test statistic and evaluated against its
corresponding null distribution obtained from randomized fits,
yielding empirical p-values that were subsequently adjusted for
multiple testing across cell types using false discovery rate (FDR)
control.


\paragraph{Predicted cell-type compositions for contrast endpoints.}
To obtain model-predicted cell-type compositions for the two endpoints
defined by a contrast, we evaluate the fitted linear model in ILR space
at the corresponding design points and then map the resulting ILR
coordinates back to the simplex. Let $\theta \in \mathbf{R}^{p}$ denote a
contrast vector defined on the rows of the coefficient matrix $B$, with
two associated design vectors $f_{\mathrm{base}}$ and $f_{\mathrm{target}}
\in \mathbf{R}^{p}$ representing the baseline and target endpoints of the
contrast, respectively (e.g.\ two factor levels being compared, with
other covariates held fixed at reference values). The predicted ILR
coordinates at each endpoint are
\begin{equation}
z_{\mathrm{base}} \;=\; f_{\mathrm{base}}^{\top} B,
\qquad
z_{\mathrm{target}} \;=\; f_{\mathrm{target}}^{\top} B,
\end{equation}
with $z_{\mathrm{base}}, z_{\mathrm{target}} \in \mathbf{R}^{D-1}$.
These ILR-space predictions are then mapped back to clr space via
\begin{equation}
\ell_{\mathrm{base}} \;=\; \Psi\, z_{\mathrm{base}},
\qquad
\ell_{\mathrm{target}} \;=\; \Psi\, z_{\mathrm{target}},
\end{equation}
yielding clr-scale log-frequency vectors in $\mathbf{R}^{D}$. Finally,
predicted cell-type proportions are obtained by applying the inverse
clr transform,
\begin{equation}
\hat{x}_{\mathrm{base},c}
\;=\;
\frac{\exp\!\left(\ell_{\mathrm{base},c}\right)}
     {\sum_{m=1}^{D} \exp\!\left(\ell_{\mathrm{base},m}\right)},
\qquad
\hat{x}_{\mathrm{target},c}
\;=\;
\frac{\exp\!\left(\ell_{\mathrm{target},c}\right)}
     {\sum_{m=1}^{D} \exp\!\left(\ell_{\mathrm{target},m}\right)},
\end{equation}
for $c = 1,\dots,D$. The resulting vectors
$\hat{x}_{\mathrm{base}}$ and $\hat{x}_{\mathrm{target}}$ are valid
compositions (non-negative and summing to one) and represent the
model-implied expected cell-type proportions at the two contrast
endpoints. These predicted compositions are uniquely defined, invariant
to the choice of reference or clr centering, and are used for
visualization and interpretation rather than for primary statistical
inference.

\paragraph{Reference set selection and reference-centered effects.}
For interpretability of cell-type loadings, we additionally define a
reference set of cell types that exhibit minimal evidence of association
with the contrast. Specifically, reference cell types are selected as
those with weak statistical evidence for differential contribution
(empirical p-value above 0.3), between one and three of them: if too few
qualify, the cell types with the smallest absolute loadings are taken,
if too many, those closest to zero are kept. Let $R \subset
\{1,\dots,D\}$ denote the resulting reference set. For visualization and
interpretation, reference-centered loadings are computed by subtracting
the mean loading over the reference set, for a cell type $c$:
\[
\tilde{\ell}_c \;=\; \ell_c - \frac{1}{|R|}\sum_{r \in R} \ell_r .
\]
This centering preserves relative differences between cell types while
anchoring effects to a set of approximately unchanged cell types. The
choice of reference affects only the presentation of loadings: the p-values are assessed as described above.


\end{document}
